\section{Critical Literature Review}
\label{sec:literature}

The empirical literature on EU~ETS carbon leakage divides by mechanism: trade-flow displacement, firm relocation and FDI, and regulated-firm performance. Table~\ref{tab:literature} summarises the identification strategy, data, main finding, and limitation of each study.

\subsection{Trade-Flow Evidence}

\citet{naegele2019} provide the most direct test of trade-flow leakage. They compute trade flows in embodied carbon and in value from Global Trade Analysis Project (GTAP) trade and input-output data over 2004 to 2011, and measure regulatory stringency as the emission cost the EU~ETS imposes on each sector, combining direct and electricity-related indirect costs net of freely allocated allowances. Estimating net-flow specifications with sector-by-country and year fixed effects, alongside bilateral specifications in the spirit of gravity models, they find no evidence that the EU~ETS has caused carbon leakage in European manufacturing. The design's strength lies in its administrative treatment measure and comprehensive sector coverage. Its limitation is aggregation: by pooling all partners and products, it cannot detect leakage concentrated in specific corridors. A separate concern is the sample period, which ends in 2011 and therefore excludes the post-2018 price acceleration.

\citet{aichele2015} estimate the carbon content of bilateral trade among 40~countries over 1995 to 2007, exploiting binding Kyoto commitment as policy variation within a theory-consistent gravity framework. They find that commitment raised a committed country's embodied carbon imports from non-committed partners by approximately 8 per cent and the emission intensity of its imports by roughly 3 per cent, though the binary indicator does not capture carbon-pricing intensity and the sample predates the post-2018 price increase.

\citet{branger2016} study cement and steel, the two most frequently cited leakage candidates, over Phases~I and~II (2005 to 2012). From an analytical model they derive an equation linking EU net imports to domestic and foreign demand and the carbon price, which they estimate with time-series methods. They find no significant effect of the carbon price on net imports of either product and conclude that competitiveness concerns in these sectors have been overstated. The analysis pools all trading partners at the EU level, so it cannot distinguish whether individual bilateral corridors behaved differently from the aggregate, and the sample ends before the Phase~III surplus was absorbed and the MSR took effect.

A common limitation spans all three trade-flow studies: none separates the import-leakage channel by partner-country regulatory stringency. If carbon-cost differentials drive leakage, the effect should be concentrated among partners with no comparable carbon pricing and pre-existing capacity in emission-intensive sectors, precisely the North African corridor this paper targets.

\subsection{Relocation and FDI Evidence}

\citet{koch2019} investigate whether EU~ETS coverage induced German multinationals to shift investment abroad, exploiting the incomplete participation requirements of the system and applying difference-in-differences with bias-corrected matching to the Bundesbank's micro database on outward foreign direct investment (FDI). The average treatment effect on non-EU FDI is close to zero and imprecisely estimated. Two margins qualify this null: less capital-intensive, geographically mobile firms increased their FDI outside the EU by 52 per cent relative to the counterfactual, and regulated firms on average expanded their number of non-EU affiliates by 28 per cent. Relocating firms are neither energy-intensive nor emission-intensive, which limits the scope for actual emissions leakage. The design is credible for the German context but pools destination countries, so it cannot identify whether investment shifted specifically toward lower-regulation partners.

\subsection{Firm-Performance Evidence}

\citet{colmer2025} estimate EU~ETS effects on French manufacturing installations over 2000 to 2012 using one-to-one nearest-neighbour Mahalanobis matching with exact sector matching and a semi-parametric matched difference-in-differences following \citet{heckman1997}. Sector is controlled through the matching procedure, not regression fixed effects. They find emissions reductions of 14 per cent in Phase~I and 16.3 per cent in Phase~II (significant at the 5\% level), emissions-intensity reductions of 17.4 per cent in Phase~II (significant at the 1\% level), and capital increases of 8.3 and 10.5 per cent (significant at the 10\% level). Value added and employment effects are not statistically significant. The pattern is consistent with abatement through capital investment rather than output contraction or relocation.

\citet{dechezlepretre2023} combine installation-level emissions data from France, the Netherlands, Norway, and the United Kingdom with firm-level performance data across EU~ETS countries in a matched difference-in-differences design with a binary treatment indicator. They find emissions reductions of approximately 10 per cent over 2005 to 2012, increases in regulated firms' revenues and fixed assets, and no statistically significant effects on profits or employment. The multi-country coverage strengthens external validity, though the binary treatment indicator does not exploit carbon-cost variation.

\citet{darcangelo2022} estimate effects on total factor productivity (TFP) for Italian manufacturing firms using structural production-function methods and find small, mostly negative effects that are heterogeneous across industries. The finding is inconsistent with a strong Porter-hypothesis interpretation. For the leakage question, the absence of large productivity losses suggests that competitive pressures from EU~ETS regulation have been modest in aggregate.

Surveys by \citet{joltreau2019} and \citet{verde2020} synthesise this evidence and reach similar conclusions: most studies find no significant competitiveness impact, attributing this to generous free allocation, though the predominance of null findings may reflect the weakness of the treatment during the low-price periods covered by most samples.

\subsection{The Gap}

Three features motivate the proposed design. First, trade-flow studies pool all EU trading partners, masking corridor-specific leakage into proximate lower-regulation partners. Second, most samples end before the EUA price exceeded EUR~25 on a sustained basis. The MSR-driven acceleration from 2018 onward remains largely unexamined. Third, free allocation attenuates the very treatment on which most studies rely. \citet{naegele2019} net freely allocated allowances out of their stringency measure, but the firm-level literature typically treats allocation as a robustness dimension rather than a component of the treatment: if free allocation fully offsets the carbon cost, the effective treatment is zero, and failure to account for this attenuation biases the estimated leakage effect toward zero.

\begin{table}[htbp]
\centering
\caption[Summary of empirical studies on EU~ETS carbon leakage]{Summary of Empirical Studies on EU~ETS Carbon Leakage}
\label{tab:literature}
\small
\renewcommand{\arraystretch}{1.3}
\begin{tabularx}{\textwidth}{>{\raggedright\arraybackslash}p{2.8cm} >{\raggedright\arraybackslash}p{2.5cm} >{\raggedright\arraybackslash}p{3.2cm} >{\raggedright\arraybackslash}p{2.8cm} >{\raggedright\arraybackslash}X}
\toprule
Author(s)/Year & Data & Identification & Main Finding & Limitation \\
\midrule
\citet{naegele2019} & Embodied-carbon trade flows, GTAP data, 2004--2011 & Net-flow and bilateral regressions; EU~ETS emission cost net of free allocation & No evidence of trade-flow leakage & Aggregation across partners and products; sample ends 2011 \\
\addlinespace
\citet{aichele2015} & Carbon content of bilateral trade, 40 countries, 1995--2007 & Kyoto commitment in a theory-consistent gravity framework & 8\% higher embodied-carbon imports from non-committed partners & Binary policy indicator; predates the EU~ETS price acceleration \\
\addlinespace
\citet{branger2016} & EU cement and steel net imports, 2005--2012 & Time-series net-import equations with demand controls & No significant carbon-price effect on net imports & EU aggregate; ends before price acceleration \\
\addlinespace
\citet{koch2019} & German multinational outward FDI, Bundesbank micro data & DiD with bias-corrected matching & No average effect; $+$52\% non-EU FDI for mobile, less capital-intensive firms & Pools destination countries; relocating firms not emission-intensive \\
\addlinespace
\citet{colmer2025} & French installations, 2000--2012 & Mahalanobis matching, semi-parametric matched DiD & Emissions $-$14\%/$-$16.3\% (5\% level); capital $+$8.3\%/$+$10.5\% (10\% level) & Single country; ends before price acceleration \\
\addlinespace
\citet{dechezlepretre2023} & Installations in 4 European countries; firm data EU-wide, 2005--2012 & Matched DiD, binary treatment & Emissions $-$10\%; revenues and assets up; no profit or employment effects & Binary treatment; does not exploit carbon-cost variation \\
\addlinespace
\citet{darcangelo2022} & Italian manufacturing firms & Structural production-function TFP estimation & Small, mostly negative, heterogeneous TFP effects & Method-dependent results; no trade-flow dimension \\
\bottomrule
\end{tabularx}
\vspace{0.3em}
\raggedright\footnotesize\textit{Source:} Own compilation based on the cited studies.
\end{table}
